\chapter{Integral lattices}\label{ch:lattices}

Everything in this book that can be computed by hand is computed in a lattice.
This chapter sets up the vocabulary --- Gram matrices, discriminants,
signatures, even and unimodular lattices, discriminant forms --- and then does
something introductory texts rarely do: it works, in complete detail, the
handful of small lattices that carry the whole later theory. The reader who
finishes the exercises will have met $\U$, $E_8$, $\U\oplus\lat{2n}$, the
K3~lattice, and the binary forms $A_2$, $\lat2\oplus\lat2$ and
$\bigl(\begin{smallmatrix}2&1\\1&4\end{smallmatrix}\bigr)$ as concrete integer
matrices, and will know exactly what it means to split one lattice off from
another.

\section{Definitions}

\begin{definition}
A \emph{lattice} is a free $\Z$-module $L$ of finite rank $r$ together with a
symmetric bilinear form $\langle\ ,\ \rangle\colon L\times L\to\Z$ that is
non-degenerate, i.e.\ $\langle x,L\rangle=0$ implies $x=0$. Choosing a basis
$(b_1,\dots,b_r)$ identifies $L$ with $\Z^r$ and the form with its \emph{Gram
matrix} $G=(\langle b_i,b_j\rangle)$, a symmetric integer matrix with
$\det G\neq0$; then $\langle x,y\rangle=x^{\mathsf T}Gy$. A change of basis by
$P\in\GL_r(\Z)$ replaces $G$ by $P^{\mathsf T}GP$, so $\det G$ is an invariant
of $L$, called the \emph{discriminant} $\disc L$. The \emph{norm} of $x$ is
$x^2:=\langle x,x\rangle$.
\end{definition}

\begin{definition}
$L$ is \emph{even} if $x^2\in2\Z$ for all $x$, equivalently if the diagonal of
$G$ is even; otherwise \emph{odd}. $L$ is \emph{unimodular} if $\disc L=\pm1$.
The \emph{signature} $(s_+,s_-)$ counts the positive and negative eigenvalues
of $G$ over $\R$ (Sylvester); $L$ is \emph{definite} if one of them is zero,
\emph{indefinite} otherwise, and \emph{hyperbolic} if $s_+=1$.
\end{definition}

\begin{definition}
For a sublattice $M\subseteq L$ of the same rank, the \emph{index} $[L:M]$ is
finite and $\disc M=[L:M]^2\disc L$; if $M$ is spanned by the columns of an
integer matrix $B$ in a basis of $L$, then $[L:M]=|\det B|$. A sublattice $M$
of any rank is \emph{primitive} if $L/M$ is torsion-free, equivalently if
$M=(M\otimes\Q)\cap L$. A vector $v$ is primitive if $\Z v$ is.
\end{definition}

\begin{definition}
The \emph{dual} of $L$ is $L^*=\{x\in L\otimes\Q:\langle x,L\rangle\subseteq\Z\}
\supseteq L$; in a basis, $L^*=G^{-1}\Z^r$. The finite group
$A_L:=L^*/L$ is the \emph{discriminant group}, of order $|\disc L|$. For $L$
even, $q_L\colon A_L\to\Q/2\Z$, $q_L(x+L)=x^2\bmod 2\Z$, is well defined and is
the \emph{discriminant form}. For $x\in L$, the \emph{divisor}
$\divv(x)$ is the positive generator of the ideal $\langle x,L\rangle\subseteq\Z$;
$x/\divv(x)\in L^*$, and $x$ is primitive if and only if $x/\divv(x)$ has order
exactly $\divv(x)$ in $A_L$.
\end{definition}

\begin{example}[the hyperbolic plane]\label{ex:U}
$\U=\Z e\oplus\Z f$ with $e^2=f^2=0$, $\langle e,f\rangle=1$, Gram matrix
$\bigl(\begin{smallmatrix}0&1\\1&0\end{smallmatrix}\bigr)$. Even, unimodular,
signature $(1,1)$. The norm form is $2xy$; $\U$ is the unique even
unimodular lattice of signature $(1,1)$ \Lstatus~\cite[Ch.~14]{Huybrechts2016}.
\end{example}

\begin{example}[the rank-one lattices]
$\lat{m}$ denotes $\Z$ with Gram matrix $(m)$. It is even iff $m$ is even, and
$A_{\lat m}\cong\Z/m$ with $q(1)=1/m\bmod2$ when $m$ is even.
\end{example}

\begin{example}[$E_8$]
The $E_8$ lattice is the unique positive-definite even unimodular lattice of
rank $8$ \Lstatus~\cite[Ch.~4]{ConwaySloane1999}; a Gram matrix is the Cartan
matrix of the root system $E_8$. Its negative $E_8(-1)$, with Gram matrix
negated, is negative-definite even unimodular. We shall never need its
entries, only its two properties: even and unimodular.
\end{example}

\begin{theorem}[\Lstatus, even unimodular lattices]\label{thm:evenunimod}
If $L$ is even and unimodular of signature $(s_+,s_-)$, then
$s_+-s_-\equiv0\pmod 8$. Conversely, for every such indefinite signature there
is exactly one even unimodular lattice, namely
$\U^{s_-}\oplus E_8^{(s_+-s_-)/8}$ when $s_+\ge s_-$, and analogously with
$E_8(-1)$ otherwise~\cite[Ch.~V]{Serre1973}, \cite[Thm.~14.1.1]{Huybrechts2016}.
\end{theorem}

\begin{definition}[the K3 lattice]
$\Lambda_{K3}:=\U^{\oplus3}\oplus E_8(-1)^{\oplus2}$, even, unimodular, of
signature $(3,19)$ and rank $22$. By Theorem~\ref{thm:evenunimod} it is the
unique even unimodular lattice of that signature.
\end{definition}

\section{Orthogonal complements and splittings}

Let $v\in L$ with $v^2\neq0$. Its orthogonal complement
$v^{\perp}=\{x\in L:\langle x,v\rangle=0\}$ is a primitive sublattice of rank
$r-1$, and $\Z v\oplus v^{\perp}$ has full rank in $L$. How far is it from all
of $L$?

\begin{proposition}[\Lstatus; elementary]\label{prop:index}
Let $v\in L$ be primitive with $v^2\neq0$, $d=\divv(v)$. Then
\[
  [L:\Z v\oplus v^{\perp}] \;=\; \frac{|v^2|}{d},
  \qquad
  \disc(v^{\perp})\cdot v^2 \;=\; \disc(L)\cdot\Bigl(\frac{v^2}{d}\Bigr)^2 .
\]
\end{proposition}

\begin{proof}
The map $L\to\Z$, $x\mapsto\langle x,v\rangle$, has image $d\Z$ and kernel
$v^{\perp}$, so $L/v^{\perp}\cong\Z$ and $L/(\Z v\oplus v^{\perp})\cong
d\Z/\langle v,v\rangle\Z$, of order $|v^2|/d$. The second identity is
$\disc(\Z v\oplus v^{\perp})=[L:\Z v\oplus v^{\perp}]^2\disc L$.
\end{proof}

In particular $\Z v\oplus v^{\perp}=L$ exactly when $\divv(v)=|v^2|$, i.e.\
when $v/v^2\in L^*$. For a $(-2)$-vector (``root'') this says: a root splits
$L$ as a direct sum iff its divisor is $2$; if its divisor is $1$ the index is
$2$.

The most useful way to \emph{exhibit} $v^{\perp}$ is not as a kernel but as a
parametrized family, together with a $2\times2$ (or smaller) Gram matrix on a
chosen basis and a \emph{frame} --- the $r\times r$ matrix whose columns are
$v$ followed by a basis of $v^{\perp}$ --- whose determinant is the index of
Proposition~\ref{prop:index}. The following worked examples are the ones that
the rest of the book depends on. They are stated for the lattice
\[
  T_N \;:=\; \U\oplus\lat{2N},\qquad
  \text{basis }(e,f,w),\quad
  \langle x,y\rangle = x_0y_1+x_1y_0+2N\,x_2y_2,\quad \disc T_N=-2N .
\]
$T_N$ is even, of signature $(2,1)$ for $N>0$, with $A_{T_N}\cong\Z/2N$ and
$q(1)=1/2N$.

\begin{leanbox}
Every equality in Examples~\ref{ex:fricke-root}--\ref{ex:disc4} below is
kernel-checked in \artifact{Agora/Geometry/MnLattice.lean} of the
repository described in Chapter~\ref{ch:formal} (declarations
\artifact{rootEF\_*}, \artifact{root7\_*}, \artifact{root7Inf\_*},
\artifact{root10Inf\_*}, \artifact{TN\_splits\_at\_rootEF},
\artifact{T7\_splits\_at\_root7}, \artifact{T7\_splits\_at\_root7Inf},
\artifact{T10\_splits\_at\_root10Inf}). The orthogonal complements are stated
there as equivalences ``$x\perp v$ iff $x$ is of the following form'', not
merely as two vectors orthogonal to $v$. Chapter~\ref{ch:formal} explains why
that distinction is the whole point.
\end{leanbox}

\begin{example}[the Fricke root, every $N$]\label{ex:fricke-root}
$v_1=e-f=(1,-1,0)$. Then $v_1^2=-2$, $\langle x,v_1\rangle=-x_0+x_1$, so
$\divv(v_1)=1$ and
\[
  x\perp v_1 \iff x=(a,a,b)=a(e+f)+b\,w .
\]
On the basis $(e+f,w)$ the form is $\operatorname{diag}(2,2N)$; so
$v_1^{\perp}\cong\lat2\oplus\lat{2N}$, of discriminant $4N$. The frame
$(v_1,e+f,w)$ has determinant $2$, and indeed $|v_1^2|/\divv=2/1=2$. Check:
$4N\cdot(-2)=-2N\cdot 2^2$.
\end{example}

\begin{example}[a second root in $T_7$]\label{ex:root7}
$v_2=(2,-4,1)=2e-4f+w$ in $T_7$. Then $v_2^2=2\cdot2\cdot(-4)+14=-2$, and
$\langle x,v_2\rangle=-4x_0+2x_1+14x_2=2(-2x_0+x_1+7x_2)$, so $\divv(v_2)=2$
and
\[
  x\perp v_2\iff x=(a,\,2a-7b,\,b)=a(1,2,0)+b(0,-7,1).
\]
On $(1,2,0),(0,-7,1)$ the Gram matrix is
$\bigl(\begin{smallmatrix}4&-7\\-7&14\end{smallmatrix}\bigr)$, of
determinant $7$; the unimodular change of basis
$\bigl(\begin{smallmatrix}2&1\\1&0\end{smallmatrix}\bigr)$ gives the reduced
basis $(2,-3,1),(1,2,0)$ with Gram matrix
$\bigl(\begin{smallmatrix}2&1\\1&4\end{smallmatrix}\bigr)$, the binary form
$x^2+xy+2y^2$ of discriminant $-7$. The frame $(v_2,(2,-3,1),(1,2,0))$ has
determinant $-1$. Since $\divv(v_2)=2=|v_2^2|$, Proposition~\ref{prop:index}
says the index is $1$: we have found an \emph{integral isometry}
\[
  \U\oplus\lat{14}\;\cong\;\lat{-2}\oplus\begin{pmatrix}2&1\\1&4\end{pmatrix}.
\]
\end{example}

\begin{example}[a class of norm $-42$ in $T_7$]\label{ex:disc3}
$v_3=(14,-14,-5)$. Then $v_3^2=-392+350=-42$,
$\langle x,v_3\rangle=14(-x_0+x_1-5x_2)$, $\divv(v_3)=14$, and
\[
  x\perp v_3\iff x=(a,\,a+5b,\,b)=a(1,1,0)+b(0,5,1).
\]
Gram matrix on that basis $\bigl(\begin{smallmatrix}2&5\\5&14\end{smallmatrix}\bigr)$,
determinant $3$; reduced basis $(1,1,0),(-2,3,1)$ with Gram matrix
$\bigl(\begin{smallmatrix}2&1\\1&2\end{smallmatrix}\bigr)$, the root lattice
$A_2$. Frame determinant $3=42/14$. Check: $3\cdot(-42)=-14\cdot3^2$.
\end{example}

\begin{example}[a class of norm $-20$ in $T_{10}$]\label{ex:disc4}
$v_4=(10,-10,-3)$ in $T_{10}$. Then $v_4^2=-200+180=-20$,
$\langle x,v_4\rangle=10(-x_0+x_1-6x_2)$, $\divv(v_4)=10$,
$x\perp v_4\iff x=(a,a+6b,b)$; reduced basis $(1,1,0),(-3,3,1)$ with Gram
matrix $\operatorname{diag}(2,2)=\lat2\oplus\lat2$; frame determinant
$2=20/10$.
\end{example}

\begin{remark}
The four examples exhibit all three behaviours of Proposition~\ref{prop:index}
in one small lattice: index $2$ (a root of divisor $1$), index $1$ (a root of
divisor $2$, a genuine direct summand), index $3$. The discriminants of the
complements --- $28,7,3$ in $T_7$ and $40,4$ in $T_{10}$ --- will reappear in
Chapter~\ref{ch:singular} as the discriminants of transcendental lattices of
K3 surfaces with Picard number $20$, and in Chapter~\ref{ch:modular} as the
discriminants of CM points on the modular curves $X_0(7)^+$ and $X_0(10)^+$.
Nothing in this chapter knows that; it is the lattice arithmetic that the
later chapters will be \emph{fed into}.
\end{remark}

\section{Embeddings, gluing, and Nikulin's theorems}

Two questions dominate the use of lattices in K3 theory: when does a lattice
$M$ embed primitively into $\Lambda_{K3}$, and when is a lattice determined by
its rank, signature and discriminant form? Both were answered by
Nikulin~\cite{Nikulin1979}; we quote the two statements used most.

\begin{theorem}[\Lstatus, Nikulin]\label{thm:nikulin-unique}
Let $L$ be an even indefinite lattice with $\rk L\ge 2+\ell(A_L)$, where
$\ell(A_L)$ is the minimal number of generators of $A_L$. Then $L$ is
determined up to isometry by its signature and discriminant form, and the
natural map $\Orth(L)\to\Orth(q_L)$ is surjective~\cite[Thm.~1.14.2]{Nikulin1979}.
\end{theorem}

\begin{theorem}[\Lstatus, Nikulin]\label{thm:nikulin-embed}
Let $M$ be an even lattice of signature $(m_+,m_-)$ and $L$ an even unimodular
lattice of signature $(\ell_+,\ell_-)$. If $\ell_+>m_+$, $\ell_->m_-$ and
$\rk L-\rk M>\ell(A_M)$, then $M$ embeds primitively into $L$, and the
orthogonal complement $M^{\perp}$ has discriminant form $-q_M$
\cite[Cor.~1.12.3, Prop.~1.6.1]{Nikulin1979}.
\end{theorem}

\begin{example}[gluing $\lat{2N}\oplus\lat{-2N}$ inside $\U$]\label{ex:glue}
In $\U$, the vectors $e+Nf$ and $e-Nf$ have norms $2N$ and $-2N$ and are
orthogonal; they span a sublattice of index $2N$ (frame determinant
$\det\bigl(\begin{smallmatrix}1&1\\N&-N\end{smallmatrix}\bigr)=-2N$). Both
are primitive. This is the elementary mechanism behind every statement of the
form ``$\lat{2N}$ and $\lat{-2N}$ are glued along $\Z/2N$ to make a unimodular
lattice'' \Kstatus~(\artifact{glue\_congruence}, \artifact{glue\_det},
\artifact{glue\_primitive}).
\end{example}

\begin{example}[$M_n$ and its complement]\label{ex:Mn}
Let $M_n:=\U\oplus E_8(-1)^2\oplus\lat{-2n}$, of signature $(1,18)$, rank
$19$, $A_{M_n}\cong\Z/2n$. By Theorem~\ref{thm:nikulin-embed} it embeds
primitively into $\Lambda_{K3}$, and any complement has rank $3$, signature
$(2,1)$ and discriminant form $-q_{M_n}$; by Theorem~\ref{thm:nikulin-unique}
(here $\rk=3\ge2+1$) that complement is unique:
\[
  (M_n)^{\perp}\;\cong\;\U\oplus\lat{2n}=T_n .
\]
This is the lattice statement underlying Dolgachev's theory of $M_n$-polarized
K3 surfaces (Chapter~\ref{ch:torelli}); it is a theorem of the literature
\Lstatus~\cite[\S7]{Dolgachev1996}. An explicit witness for $n=7$ inside
$\U^3$ --- integer matrices $B,C$ with $B^{\mathsf T}\U^3B=T_7$,
$C^{\mathsf T}\U^3C=\U\oplus\lat{-14}$, $B^{\mathsf T}\U^3C=0$, the $E_8(-1)^2$
summand splitting off because it is unimodular --- is kernel-checked
\Kstatus~(\artifact{Agora/Geometry/Embedding.lean}). What is \emph{not}
kernel-checked there is primitivity of the images and the identification of
each as the other's complement; that step is Theorem~\ref{thm:nikulin-embed}.
\end{example}

\section{Isometries and reflections}

For a vector $v$ with $v^2\neq0$ and $2\langle x,v\rangle/v^2\in\Z$ for all
$x$ --- automatic for a root, $v^2=-2$ --- the \emph{reflection}
\[
  s_v(x)=x-\frac{2\langle x,v\rangle}{v^2}\,v
\]
is an integral isometry of $L$ fixing $v^{\perp}$ pointwise and sending $v$ to
$-v$. Reflections in roots generate the Weyl group of $L$, and the walls
$v^{\perp}\otimes\R$ cut the positive cone into chambers; for a K3 surface these
chambers are where the ample classes live (Chapter~\ref{ch:torelli}).

\begin{example}[the swap]\label{ex:swap}
In $T_N$, the exchange $\sigma\colon e\leftrightarrow f$, $w\mapsto w$, is an
isometry, an involution, of determinant $-1$; it is the reflection in the
Fricke root $e-f$ \Kstatus~(\artifact{swap\_isometry},
\artifact{swap\_involution}). Its fixed lattice is $v_1^{\perp}=\lat2\oplus\lat{2N}$
of Example~\ref{ex:fricke-root}. In Chapter~\ref{ch:modular} this same matrix
will be the Fricke involution $\tau\mapsto-1/(N\tau)$.
\end{example}

\section{Exercises}

\begin{exercise}
Show that $\Z/2N$ with $q(1)=1/2N$ and $\Z/2N$ with $q(1)=-1/2N$ are the
discriminant forms of $\lat{2N}$ and $\lat{-2N}$, and that they are
isomorphic as finite quadratic forms iff $-1$ is a square modulo $4N$ (for
even $N$) resp.\ modulo $2N$ (for odd $N$)... or rather: determine exactly for
which $N\le 12$ they are isomorphic. (Hint: Appendix~\ref{app:hints}.)
\end{exercise}

\begin{exercise}[$\star$]\label{exr:disc-identity}
Prove the identity $\disc(v^{\perp})\cdot v^2=\disc(L)\cdot(v^2/\divv v)^2$
of Proposition~\ref{prop:index} directly from a frame, i.e.\ from
$\det(F^{\mathsf T}GF)=\det(F)^2\det G$ applied to the frame $F=(v\mid B)$
with $B$ a basis of $v^{\perp}$.
\end{exercise}

\begin{exercise}
In $T_7$, find all roots $v$ (norm $-2$) with $|v_0|,|v_1|,|v_2|\le 5$, compute
$\divv(v)$ for each, and sort them into the two classes ``index $2$'' and
``index $1$''. Which reduced binary forms occur as $v^{\perp}$?
\end{exercise}

\begin{exercise}[$\star$]
Show that in $T_N$ every vector of the form $(N,-N,-k)$ has
$\langle x,v\rangle=N(-x_0+x_1-2kx_2)$, hence $\divv(v)=N$ when $\gcd(N,2k)=N$
divides appropriately, and compute $v^2$ and $\disc(v^{\perp})$ as functions
of $N,k$. Recover Examples~\ref{ex:disc3} ($N=7,k=5$) and \ref{ex:disc4}
($N=10,k=3$).
\end{exercise}

\begin{exercise}
Prove that $\U\oplus\lat{2N}$ and the lattice with Gram matrix
$\bigl(\begin{smallmatrix}0&0&-1\\0&2N&0\\-1&0&0\end{smallmatrix}\bigr)$
are isometric, by an explicit $P\in\GL_3(\Z)$. (Both conventions occur in the
literature; Chapter~\ref{ch:frontier} records a case where the choice caused
confusion between two teams.)
\end{exercise}

\begin{exercise}
The lattice $G_0(N)$ with Gram matrix $\bigl(\begin{smallmatrix}0&0&-2\\0&1&0\\-2&0&0\end{smallmatrix}\bigr)\cdot N$-adjusted as in
Chapter~\ref{ch:frontier}, of discriminant $-4N^2$, carries an action of
$\Gamma_0(N)$ by symmetric squares. Show that no lattice of discriminant
$-4N^2$ can be isometric to $T_N$ unless $N\in\{0,\tfrac12\}$. (This trivial
observation is the content of a kernel-checked theorem,
\artifact{no\_isometry\_G0N\_TN}, recorded because the two lattices were once
conflated under one name.)
\end{exercise}

\section*{Chapter note}

Definitions and Theorems~\ref{thm:evenunimod}, \ref{thm:nikulin-unique},
\ref{thm:nikulin-embed} follow Huybrechts~\cite[Ch.~14]{Huybrechts2016} and
Nikulin~\cite{Nikulin1979}; statements were checked against those sources.
Proposition~\ref{prop:index} is folklore; the proof given is complete. The
numerical content of Examples~\ref{ex:fricke-root}--\ref{ex:disc4} is
kernel-checked as described; the reader is encouraged to redo each by hand,
which takes ten minutes and is the best possible introduction to the subject.
